\newpage
\begin{problem}[Warmup]{(a)}
    Give examples of a commutative ring with unity and an ideal satisfying the following (or argue that no such example exists):
    \begin{enumerate}
        \item[(a)] the ideal is prime but not maximal.
        \item[(b)] the ideal is maximal but not prime.
        \item[(c)] the ideal is prime and the ring has zero divisors.
    \end{enumerate}
\end{problem}
\unsolved

\begin{problem}[Warmup]{(b)}
    Explain why $\langle x^2 - 2 \rangle$ is a maximal ideal of $\mathbb{Q}[x]$, but not of $\mathbb{R}[x]$.
\end{problem}
\unsolved

\begin{problem}[Warmup]{(c)}
    Explain why $\langle x - 2 \rangle$ is a prime ideal of $\mathbb{Z}[x]$ but is not maximal.
\end{problem}
\unsolved

\begin{problem}[Warmup]{(d)}
    Give the quickest possible proof that a maximal ideal is prime.
\end{problem}
\unsolved

\begin{problem}[Warmup]{(e)}
    What are the zero divisors of $\mathbb{Z} \times \mathbb{Z}$? What about of $\mathbb{Z}_4 \times \mathbb{Z}$?
\end{problem}
\unsolved

\begin{problem}[Warmup]{(f)}
    Let $R$ be a commutative ring with unity. Under what circumstances is $\{0\}$ a prime ideal? A maximal ideal?
\end{problem}
\unsolved

\begin{problem}[Warmup]{(g)}
    Let $R$ be a commutative ring with unity. Why isn't the ideal $R$ maximal or prime?
\end{problem}
\unsolved

\begin{problem}{1}
    Let $R$ be a commutative ring with unity, $I$ an ideal of $R$, and $a \in R$. Define
    \[\langle I, a \rangle = \{k + ab : k \in I, b \in R\}.\]
    Prove that $\langle I, a \rangle$ is an ideal. Then show that
    $I \subseteq \langle I, a \rangle$ and $a \in \langle I, a \rangle$. Furthermore, show
    that $\langle I, a \rangle$ is the smallest ideal of $R$ that contains both $I$ and
    $a$. (This result is needed in the proof of Theorem 20.1.)
\end{problem}
\unsolved

\begin{problem}{2}
    Prove Theorem 20.3. That is, suppose the $R$ is a commutative ring,
    and $0 \ne a \in R$. Prove that $\langle a \rangle$
    is a prime ideal if and only if $a$ is a prime element.
\end{problem}
\solved

\begin{proof}
    Suppose $\langle a \rangle$ is a prime ideal.
    Now, suppose $a \mid xy$.
    Then $at = xy$ for some $t \in R$, so $xy \in \langle a \rangle$.
    We must show either $a \mid x$ or $a \mid y$.
    Since $\langle a \rangle$ is a prime ideal, we have $x \in \langle a \rangle$ or $y \in \langle a \rangle$.
    Suppose, WLOG, $x \in \langle a \rangle$.
    Then $x = ak$ for some $k \in R$, so $a \mid x$, as required.

    Conversely, suppose $a$ is a prime element.
    Let $xy \in \langle a \rangle$.
    Then $xy = ak$ for some $k \in R$.
    We must show either $x \in \langle a \rangle$ or $y \in \langle a \rangle$.
    So $a \mid xy$, and, since $a$ is prime, $a \mid x$ or $a \mid y$.
    Suppose, WLOG, $a \mid x$.
    So $x = af$ for some $f \in R$.
    It follows that $x \in \langle a \rangle$,
        so $\langle a \rangle$ is a prime ideal, as required.
\end{proof}

\begin{problem}{3}
    Let $p$ be a prime integer in $\mathbb{Z}$. Prove directly that $\langle p \rangle$ is
    a maximal ideal.
\end{problem}
\unsolved

\begin{problem}{4}
    Find all prime ideals of $\mathbb{Z}$. Find all maximal ideals of $\mathbb{Z}$.
\end{problem}
\solved

\textbf{Solution: } The ideals of the form $\langle p \rangle$
(all ideals are of this form since $\mathbb{Z}$ is a PID) where $p$ is prime (which corresponds to irreducibility in $\mathbb{Z}$)
are maximal, and these are all of the maximal ideals.
The prime ideals are exactly the maximal ideals together with $\langle 0 \rangle$.

\begin{problem}{5}
    Find all prime ideals of $\mathbb{Z}_{20}$. Find all maximal ideals of
    $\mathbb{Z}_{20}$.
\end{problem}
\unsolved

\begin{problem}{6}
    Find all prime ideals of $\mathbb{Z} \times \mathbb{Z}$. Find all maximal ideals of
    $\mathbb{Z} \times \mathbb{Z}$.
\end{problem}
\unsolved

\begin{problem}{7}
    Find all prime ideals of $\mathbb{Z}_2 \times \mathbb{Z}_3$; find all maximal ideals of
    $\mathbb{Z}_2 \times \mathbb{Z}_3$. Now do the same for $\mathbb{Z}_2 \times \mathbb{Z}_4$.
\end{problem}
\unsolved

\begin{problem}{8}
    Find all prime ideals of $\mathbb{Q}$ and $\mathbb{Q} \times \mathbb{Q}$; do likewise for all maximal ideals.
\end{problem}
\unsolved

\newpage

\begin{problem}{9}
    Consider $\mathbb{Q}[x, y]$, the set of all polynomials with coefficients from
    $\mathbb{Q}$, in the two indeterminates $x$ and $y$. In Exercise 12.12, you considered
    this ring and showed that it is not a PID.
    \begin{enumerate}
        \item[(a)] Consider the function $\psi : \mathbb{Q}[x, y] \to \mathbb{Q}$ defined by
        $\psi(f) = f(0, 0)$. Why is this a ring homomorphism? Why is the kernel of $\psi$
        a maximal ideal? Determine explicitly a nice description of the kernel.
        \item[(b)] Consider the function $\rho : \mathbb{Q}[x, y] \to \mathbb{Q}[x]$ defined
        by $\rho(f) = f(x, 0)$. Why is this a ring homomorphism? Why is the kernel of
        $\rho$ a prime ideal? Determine explicitly a nice description of the kernel. Find a
        larger proper ideal, thus showing concretely that this ideal is not maximal.
        \item[(c)] Consider the function $\varphi : \mathbb{Q}[x, y] \to \mathbb{Q}[x]$
        defined by $\varphi(f) = f(x, x)$. Prove that this is a ring homomorphism. Why is
        the kernel of $\varphi$ prime but not maximal? Determine explicitly a nice
        description of the kernel. Find a larger proper ideal, thus showing concretely
        that this ideal is not maximal.
    \end{enumerate}
\end{problem}
\unsolved

\begin{problem}{10}
    Consider the function $\varphi : \mathbb{Z} \to \mathbb{Z} \times \mathbb{Z}$ defined
    by $\varphi(n) = (n, n)$. Check that this is a homomorphism. Now,
    $\ker(\varphi) = \{0\}$. Argue directly from the definition that $\{0\}$ is a prime
    ideal of $\mathbb{Z}$. Now $\mathbb{Z} \times \mathbb{Z}$ is obviously not a domain;
    why does this not contradict the theorem in the text characterizing prime ideals?
\end{problem}
\unsolved

\begin{problem}{11}
    Consider the function $\varphi : \mathbb{Q}[x] \to \mathbb{R}$ defined by
    $\varphi(f) = f(\sqrt{2})$. Check that this is a homomorphism.
    \begin{enumerate}
        \item[(a)] Describe the elements of the ring $\varphi(\mathbb{Q}[x])$.
        \item[(b)] This ring is clearly a domain (because it is a subring of $\mathbb{R}$).
        Is it a field?
        \item[(c)] What is the kernel of this homomorphism? Is it prime or maximal?
        \item[(d)] How could we reach the conclusion for part c, using an argument from
        Chapter 13?
    \end{enumerate}
\end{problem}
\unsolved

\begin{problem}{12}
    Consider the function 
    \[\rho : \mathbb{Z}[x] \longrightarrow \mathbb{Z}_3 [x]\]
    defined by 
    \[\rho(a_0 + a_1 x + \cdots + a_n x^n) = [a_0] + [a_1]x + \cdots + [a_n]x^n.\]
    \begin{enumerate}
        \item Prove that this is an onto ring homomorphism.
        \item Why is the kernel of $\rho$ a prime ideal? Give an explicit description of its kernel (it is a principal ideal).
    \end{enumerate}
\end{problem}
\solved

\begin{proof}
    Let $f = \sum a_i x^i$ and
    $g = \sum b_i x^i$ in $\mathbb{Z}[x]$, the $i$th coefficient of $f + g$
    is $a_i + b_i$, and $[a_i + b_i] = [a_i] + [b_i]$, so
    $\rho(f + g) = \rho(f) + \rho(g)$.
    Similarly, the $k$th coefficient of $fg$ is $\sum_{i+j=k} a_i b_j$, and
    $\left[\sum_{i+j=k} a_i b_j\right] = \sum_{i+j=k} [a_i][b_j]$,
    so $\rho(fg) = \rho(f)\rho(g)$.
    Finally, $\rho(1) = [1]$.

    Now, let $[b_0] + [b_1]x + \cdots + [b_n]x^n \in \mathbb{Z}_3[x]$,
    where $b_i \in \mathbb{Z}$.
    Notice, $b_0 + b_1 x + \cdots + b_n x^n \in \mathbb{Z}[x]$ and
    $\rho(b_0 + b_1 x + \cdots + b_n x^n) = [b_0] + [b_1]x + \cdots + [b_n]x^n$.
    Thus $\rho$ is onto.

    Now, suppose $\rho(a_0 + a_1 x + \cdots + a_n x^n) = [a_0] + [a_1]x + \cdots + [a_n]x^n = 0$.
    It follows that $a_i \equiv 0 \pmod{3}$ for each $i$, so every coefficient
    is divisible by $3$. Conversely, any polynomial with all coefficients
    divisible by $3$ maps to $0$.
    Thus $\ker \rho = \langle 3 \rangle$, the principal ideal generated by $3$
    in $\mathbb{Z}[x]$.
    Since $\rho$ is onto, the First Isomorphism Theorem gives
    $\mathbb{Z}[x]/\ker \rho \cong \mathbb{Z}_3[x]$, which is an integral
    domain because $\mathbb{Z}_3$ is. Therefore $\ker \rho$ is a prime ideal.
\end{proof}

\begin{problem}{13}
    Consider the ideal $\langle x^2 - 2 \rangle$ in $\mathbb{Q}[x]$, discussed in
    Example 20.5. Prove that this is a maximal ideal, without using either Chapter 13 or
    Theorem 20.1.
\end{problem}
\unsolved

\begin{problem}{14}
    Consider the function $\varphi : \mathbb{Z}[x] \to \mathbb{Q}$ defined by
    $\varphi(f) = f(-1/2)$.
    \begin{enumerate}
        \item[(a)] What is the kernel of $\varphi$?
        \item[(b)] We know that $\mathbb{Q}$ is a field. Does this mean that
        $\ker(\varphi)$ is a maximal ideal? Be careful!
        \item[(c)] Give a nice description of the elements $\varphi(\mathbb{Z}[x])$.
    \end{enumerate}
\end{problem}
\unsolved

\begin{problem}{15}
    Consider the homomorphism discussed in Exercise 16.30 and Exercise 19.11. What ideal
    do we now know is maximal? Prove that this ideal is maximal directly, from the
    definition.
\end{problem}
\unsolved

\begin{problem}{16}
    Consider the homomorphism discussed in Exercise 16.10 and Exercise 19.9. What ideal
    do we now know is maximal? Prove that this ideal is maximal directly, from the
    definition.
\end{problem}
\unsolved

\begin{problem}{17}
    Let $R$ be a commutative ring with unity, and $I$ a prime ideal of $R$.
    NOte that $I[x]$ is a subring of $R[x]$. Prove that $I[x]$ is a prime ideal of $R[x]$.
    Now suppose that $I$ is a maximal ideal; show by example that $I[x]$ need not be maximal.
\end{problem}
\unsolved

\begin{problem}{18}
    Let $R$ and $S$ be commutative rings. Prove that if the direct product $R \times S$ is
    a domain, then exactly one of $R$ and $S$ is the zero ring.
\end{problem}
\unsolved

\begin{problem}{19}
    It turns out that in a commutative ring with unity, every proper ideal is a subset of
    a maximal ideal. (The proof of this theorem requires more sophisticated set theory
    than we wish to enquire into here.) Use this theorem to establish the following:
    \begin{enumerate}
        \item[(a)] Prove that every commutative ring with unity has a field as a
        homomorphic image.
        \item[(b)] If the commutative ring $R$ with unity has a unique maximal ideal $M$,
        then $M$ consists of exactly the set of non-units.
    \end{enumerate}
\end{problem}
\unsolved

\begin{problem}{20}
    Let $P$ be an ideal of a commutative ring with unity. Prove that $P$ is a prime ideal
    if and only if, whenever $P \supseteq I \cdot J$, where $I$ and $J$ are ideals, then
    $P \supseteq I$ or $P \supseteq J$. (For the definition of the product of two ideals,
    see Exercise 11.15.)
\end{problem}
\unsolved

\begin{problem}{21}
    Prove that in a PID, a non-zero proper ideal is prime if and only if it is maximal.
\end{problem}
\unsolved

\begin{problem}{22}
    Prove that in a finite commutative ring with unity, a proper ideal is prime if and
    only if it is maximal.
\end{problem}
\unsolved

\begin{problem}{23}
    Let $R$ be a commutative ring with unity, for which every element is an idempotent;
    such a ring is called \textbf{Boolean}. (See Exercise 7.25 for more about
    idempotents.) Prove that in $R$, a proper ideal is prime if and only if it is maximal.
\end{problem}
\unsolved
